\chapter{Introduction}\label{ch:introduction}

\section{Background and Motivation}\label{sec:intro-background}

Video retrieval is a central task in multimedia information systems, supporting content search, recommendation, personalization, and large-scale video discovery. Given a query, a retrieval system identifies relevant videos from a candidate collection and provides a compact set of candidates for subsequent ranking or decision-making~\cite{zhu2023deep}. Early research formulated this as text-to-video matching over short clips, evaluated by the alignment between text queries and corresponding clips. The field has progressed from adapting Contrastive Language-Image Pretraining (CLIP)~\cite{radford2021learning} to video through temporal modules~\cite{luo2022clip4clip} to employing Multimodal Large Language Models (MLLMs) that offer superior language understanding and visual generalization~\cite{kong2025modality}. These models are typically trained on massive datasets with noisy text annotations~\cite{bain2021frozen} and evaluated on standard benchmarks such as MSR-VTT~\cite{xu2016msr}, achieving strong results for coarse-grained text-to-video retrieval. In parallel, retrieval-and-ranking architectures~\cite{cheng2016wide,covington2016deep} have become the standard deployment paradigm for video search and recommendation systems. Contemporary video retrieval, however, extends well beyond this controlled regime. Retrieval systems increasingly encounter heterogeneous query forms, richer semantic conditions, longer and more informationally dense videos, and candidate collections whose size precludes exhaustive scoring. These developments shift video retrieval from a fixed benchmark matching problem to a broader methodological question: how can retrieval scale beyond the controlled conditions under which it was originally developed?

Building a video retrieval system involves three decisions in sequence, each presenting a distinct scaling challenge. First, the semantic capacity of the retrieval model must scale to the diversity of queries and relevance patterns encountered in practice. Second, the observation strategy through which visual evidence is provided to the model must scale to videos of varying length and information density. Third, the learned retrieval signal must scale to large candidate collections while satisfying the latency and throughput constraints of real serving systems. We term these three challenges semantic scaling, observational scaling, and search scaling. They follow the information flow of the retrieval pipeline, where each stage depends on the output of the preceding one. Yet each presents a distinct scaling bottleneck that is not eliminated by optimizing an upstream stage alone. A better representation does not automatically determine how videos should be observed, and a better observation strategy does not automatically make retrieval signals searchable over billion-scale collections. Addressing scalability in video retrieval therefore requires understanding the specific bottleneck at each stage and designing solutions targeted at its structure.

Semantic scaling asks whether the retrieval representation can scale to the full breadth of relevance conditions encountered in practice: textual, visual, or composed queries, global scene-level semantics or fine-grained event details, temporal dynamics, spatial relations, and partial relevance within long videos. However, the evaluation landscape is fragmented. Recent benchmarks have advanced specific aspects such as spatiotemporal grounding~\cite{xu2025carebench}, composed retrieval~\cite{ventura2024covr}, and long-context video understanding~\cite{cai2025lovr}, but models are optimized and tested on isolated benchmarks, preventing a clear understanding of their comprehensive capabilities. This fragmentation obscures a deeper problem. Performance on conventional benchmarks correlates poorly with holistic retrieval capability, and scaling model size alone does not resolve these deficiencies. Image-only trained models~\cite{Zhang_2025_CVPR} can perform surprisingly well on video retrieval tasks, and compact models can outperform much larger ones on specific retrieval scenarios~\cite{wang2024internvideo2}, suggesting that expanding model parameters and data volume does not by itself yield broader retrieval capability. Two additional gaps compound the problem: there is no systematic diagnostic framework for identifying which capabilities a model is missing, and existing data synthesis approaches~\cite{chen2023vast,chai2024auroracap} produce resources with uneven quality and distribution. Without knowing which capabilities to scale toward, representation learning expands in directions that may not address the most consequential deficiencies.

Observational scaling asks whether the observation strategy can scale to videos of varying length and information density. Before a video can be processed by an embedding model, it must be converted into a discrete sequence of visual inputs through spatiotemporal sampling. Since encoding every frame at full resolution is computationally prohibitive~\cite{hu2022mmsampler}, this sampling step fundamentally determines the visual evidence available to the representation~\cite{guo2025gve}. Too few frames or too low a resolution can cause the target event to be missed entirely, constituting a structural information bottleneck that cannot be addressed by representation learning alone. Yet most existing retrieval research operates in settings where this constraint is invisible. Short-clip datasets such as MSR-VTT~\cite{xu2016msr} and DiDeMo~\cite{anne2017localizing} are brief enough that sparse sampling captures most of the video content. Real-world videos, however, are increasingly long and informationally dense~\cite{henschel2025streamingt2v,weng2024longvlm}, and user queries may describe anything from the overall theme to a brief, specific moment, making retrieval inherently multi-granular~\cite{xiao2022video,chen2023joint}. How sampling density interacts with query granularity is a central question for long-video retrieval~\cite{wang2025cross}, yet this relationship remains underexplored because no existing benchmark provides the required testbed. Short-clip benchmarks are too brief for sampling density to matter, long-video benchmarks~\cite{cai2025lovr} evaluate at a single query granularity, and benchmarks with fine-grained temporal annotations feature videos that are relatively short. Without a setting that combines long-form visual evidence with queries at multiple granularity levels, the scaling behavior of retrieval with respect to observation remains poorly understood. Recent work has begun to study scaling laws for frame count and caption density on short clips~\cite{zhu2025dynamic}, but the interaction between spatiotemporal sampling and retrieval granularity on long-form video has not been investigated.

Search scaling asks whether learned retrieval signals remain applicable under the computational constraints of large candidate collections. Video retrieval and recommendation systems in production operate over corpora that may contain hundreds of millions to billions of items~\cite{wu2022survey,zhu2018learning}, where exhaustive scoring is incompatible with online latency requirements~\cite{liu2021jizhi}. Graph-based approximate nearest neighbor search methods such as HNSW~\cite{malkov2018efficient} have become the dominant approach, but a structural conflict exists between learned retrieval metrics and graph-based search. Neural relevance metrics often violate the triangle inequality~\cite{chen2022approximate}, the property that ensures optimal pathfinding during graph traversal, forcing the exploration of redundant edges and causing greedy traversal to become trapped in local regions of the graph. Beyond metric-search incompatibility, high concurrency introduces hardware execution bottlenecks. Divergent queries trigger different search paths, producing irregular neural computations that lead to frequent GPU kernel launches with fixed latency overhead dominating runtime~\cite{el2016klap,liu2022lion}. These challenges are tightly coupled: the learned metric affects graph traversal behavior, the traversal algorithm determines computational patterns, and the hardware execution strategy constrains achievable throughput. Optimizing any single component in isolation does not resolve the overall system bottleneck.

To address these three challenges, the studies in this thesis develop targeted solutions at each stage of the retrieval pipeline. The first study establishes a unified evaluation protocol across sixteen public datasets to diagnose capability gaps, develops a synthesis pipeline (V-SynFlow) that produces over 1.55 million multi-task video-text pairs targeting the identified deficiencies, and applies a curriculum learning strategy (Modality Pyramid) that exploits inherent task hierarchies for progressive training of a general video embedder. The second study introduces the Multi-Granular Long Video Retrieval (MG-LVR) benchmark, conducts a comprehensive spatiotemporal scaling analysis across twenty-three models that reveals a consistent asymmetry between spatial and temporal scaling behavior, and proposes a dense-to-sparse self-distillation framework~\cite{hinton2015distilling} (VEterfall) that transfers retrieval capability from high-density observation to efficient sparse inference. The third study co-designs a learned metric with approximate triangle inequality for graph search compatibility, a parallel graph search algorithm with breadth-depth-balanced traversal, and GPU-accelerated execution with adaptive aggregation, validated through deployment in a large-scale video recommendation system serving hundreds of millions of users.

Together, these studies form a coherent trajectory along the retrieval pipeline, from semantic representation to visual observation and large-scale search execution. Across the three stages studied in this thesis, the results show that scaling model capacity, observation density, or metric accuracy alone is insufficient for broader retrieval generalization unless it is aligned with the structure of the corresponding bottleneck. Scalable video retrieval therefore requires structure-aware scaling at each stage. Within this thesis, generalizability serves as an empirical criterion for assessing whether such scaling has been performed effectively. By addressing semantic, observational, and search scaling in sequence, this thesis develops a methodology for video retrieval that extends from representation learning through observation-aware embedding generation to deployment-scale search, providing a path toward retrieval systems that are both scalable and generalizable.

\section{Organization of the Thesis}\label{sec:intro-organization}

The remainder of this thesis is organized as follows.

Chapter~\ref{ch:related-work} reviews the background and related literature, covering fundamental concepts and prior work on video embedding models, benchmarks and evaluation protocols, and large-scale retrieval.
Chapter~\ref{ch:universal-video-retrieval} addresses semantic scaling through capability-guided evaluation, targeted data synthesis, and curriculum learning, producing the General Video Embedder (GVE). This study establishes the representation foundation that the subsequent studies build upon.
Chapter~\ref{ch:spatiotemporal-scaling} addresses observational scaling by introducing the Multi-Granular Long Video Retrieval (MG-LVR) benchmark, analyzing spatiotemporal scaling behavior, and developing the Video Embedding Waterfall (VEterfall) framework for efficient inference. This study extends the investigation from what the model should represent to how visual evidence is provided to the model.
Chapter~\ref{ch:scale-efficient-retrieval} addresses search scaling through co-designed learned metrics, parallel graph search (HiPANNS), and GPU-accelerated execution, validated through deployment in a large-scale video recommendation system. This study addresses the final stage of the retrieval pipeline.
Chapter~\ref{ch:conclusion} summarizes the main contributions and discusses future directions. Chapter~\ref{ch:publications} lists the publications associated with this thesis.
